% ---------------------------------------------------------------------
\section{Önerilen Yöntem}\label{sec:method}
% ---------------------------------------------------------------------

\Fig{fig:architecture} üzerinde görüldüğü gibi \method{}, her adımda art arda çalışan üç parçadan oluşur.
Dikkat tabanlı politika, bir şarj ünitesi belirteçleri kümesini şarj ünitesi başına bir ham eyleme eşler; güvenlik katmanı ise ham eylemleri araç başına sınırlara ve trafo limitine uyan komut güçlerine dönüştürür.
Eğitim sırasında bir politika gradyanı yöntemi, güvenlik katmanını ortamın bir parçası olarak ele alarak politikayı gerçekleşen kâra göre günceller; aynı katman kullanım sırasında da değişmeden çalışır.

\begin{figure}[t]
\centering
\figorbox{figures/fig_architecture.pdf}{\linewidth}{6cm}
\caption{\method{} yönteminin mimarisi ve kontrol döngüsü. Düz oklar eğitimde ve kullanımda işletilen kontrol döngüsünü oluşturur; kesikli oklar yalnızca eğitimde kullanılır. Güvenlik katmanı (gölgeli), trafo limitini koruyan ve her EV alt sınırını aldığı sürece EV'nin gereksinimini erişilebilir tutan kesin bir izdüşümdür (Teorem~\ref{thm:invariance}).}
\label{fig:architecture}
\end{figure}

% ---------------------------------------------------------------------
\subsection{Belirteç gösterimi ve dikkat tabanlı politika}\label{sec:policy}

Gözlem, $N$ adet şarj ünitesi belirteci $\xi_i(t)\in\R^{12}$ ve bir küresel belirteç $\xi_0(t)\in\R^{38}$ içeren bir kümedir; hiçbir öznitelik bir şarj ünitesinin indisini ya da konumunu kodlamaz.
Bir şarj ünitesi belirteci doluluk, enerji ve kapasite verilerini, kalan gerekli enerjiyi ve adım sayısını, \Sec{sec:bounds} içinde tanımlanan esneklik payını, güç limitlerini ve $l_i$ ile $u_i$ sınırlarını içerir.
Küresel belirteç sonraki $H_f=8$ adımın fiyatlarını, zaman özniteliklerini, mevcut limitleri, yük ve PV tahminlerini ve dört istasyon oranını içerir.
Tüm öznitelikler $N$ ile büyümeyecek biçimde normalize edilir; Tablo~\ref{S-tab:features} bunları listelemektedir.

Her belirteç, öğrenilen iki tür gömmesinden biriyle birlikte doğrusal olarak gömülür: şarj üniteleri için $e_i=W_c\xi_i+b_c+v_c$, küresel belirteç için $e_0=W_g\xi_0+b_g+v_g$.
$[e_0;e_1;\dots;e_N]$ dizisi, genişliği $d_m=128$ olan, $n_H=4$ başlıklı ve konum kodlaması içermeyen $n_L=3$ adet ön-LayerNorm'lu Transformer kodlayıcı katmanından geçer \cite{vaswani2017attention,ba2016layer,xiong2020layer}.
Bir anahtar dolgu maskesi boş şarj ünitelerini dikkat anahtarları arasından çıkarır; küresel anahtar ise hiçbir zaman maskelenmez.
Ardından ortak bir aktör başlığı, her şarj ünitesi çıktısı $h_i$'yi Gauss dağılımlı bir ham eylemin $\mu_{\theta,i}$ ortalamasına ve $\sigma_{\theta,i}$ standart sapmasına eşler; burada $\theta$ aktör parametrelerini gösterir.
Politika $\pi_\theta(a\mid s)=\prod_{i\in\mathcal N(t)}\mathrm{Normal}\big(a_i;\mu_{\theta,i}(s),\sigma_{\theta,i}^2(s)\big)$ biçimindedir.
$\omega$ parametreli ayrı bir kritik kodlayıcısı, maskelenmemiş tüm belirteç çıktılarını öğrenilen tek bir dikkat sorgusuyla \cite{lee2019set} havuzlayarak $V_\omega(s)$ değerini üretir.

Bu yapı permütasyona eşdeğişkendir (permutation-equivariant); maskelenmiş belirteçler maskelenmemiş belirteçlerin çıktılarını etkilemez ve parametre sayısı $N$'ye bağlı değildir (Lemma~\ref{S-lem:equivariance}).
Şarj ünitesi sayısı ne olursa olsun aktörün 610\,370, kritiğin ise 676\,481 parametresi vardır.
Dolayısıyla $N\in\{10,20\}$ ile eğitilen ağırlıklar $N=200$ için çalıştırılabilir; ancak bu boyutta performans garanti edilmez, çünkü 200 belirteç üzerindeki dikkat, eğitimde görülenden daha fazla anahtar üzerinden ortalama alır.

% ---------------------------------------------------------------------
\subsection{Erişilebilir enerjiden türetilen araç başına sınırlar}\label{sec:bounds}

Güvenlik katmanı, bağlı her EV için ayrılış gereksinimini erişilebilir tutan en küçük güce ihtiyaç duyar; karışıklığa yol açmadığı yerlerde $i$ indisi yazılmamaktadır.
Bir adım boyunca tam güçte şarj, $E$ enerjisini $\varphi(E)=E+\Delta t\,\eta^{\sm c}\,\bar P(E)$ değerine eşler.
\Sec{sec:invariance} içindeki A1 varsayımı, $\varphi$'nin sürekli ve azalmayan olmasını, yani bir adıma daha fazla enerjiyle başlamanın tam güçte şarjdan sonra hiçbir zaman daha az enerji bırakmamasını gerektirir.
Gereksinime tam güçte hâlâ erişilebilen en küçük enerjiler olan geriye doğru erişilebilirlik eşikleri şunlardır:
\begin{equation}\label{eq:thresholds}
B(0)=\Ereq,\qquad B(k)=\min\big\{E\in[E^{\min},E^{\sm{cap}}]:\ \varphi(E)\ge B(k-1)\big\},\quad k\ge1 .
\end{equation}
A1 altında bir EV, $t$ adımında bireysel olarak olurludur (yani kabul edilebilir bir kontrol \Eq{eq:requirement} koşulunu sağlar) ancak ve ancak $E(t)\ge B(\tau(t))$ ise (Lemma~\ref{S-lem:reach}).
A1 olmadan bu nitelendirme geçerliliğini yitirir, çünkü şimdi daha az şarj etmek daha sonra daha fazla erişilebilir enerji bırakabilir.
EV2Gym'in üstel şarj gücü azalması için Lemma~\ref{S-lem:taper}, A1'i her adım uzunluğu için ispatlamaktadır.

$E=E_i(t)$ enerjisine ve $\tau=\tau_i(t)$ kalan eyleme sahip bağlı bir EV için güvenlik katmanının sınırları şunlardır:
\begin{equation}\label{eq:bounds}
u_i(t)=\bar P(E),\qquad
l_i(t)=\begin{cases}
\min\Big\{u_i(t),\ \max\Big\{\pmn(E),\ g^{-1}\Big(\dfrac{\Bh(\tau-1)-E}{\Delta t}\Big)\Big\}\Big\}, & E\ge\Bh(\tau),\\[2mm]
\bar P(E), & E<\Bh(\tau),
\end{cases}
\end{equation}
burada $\Bh$ saklanan bir eşik tablosu, dıştaki minimum ise sayısal bir korumadır.
Birinci durumda $l_i$, ardından gelen $\tau-1$ tam güç adımının $\Ereq$ değerine hâlâ ulaştığı kabul edilebilir en küçük güçtür; bireysel olarak olursuz bir EV ise en yüksek güçte şarj edilir.
Negatif bir alt sınır, pay kaldığı sürece V2G deşarjına izin verir ve \Fig{fig:concepts}(a) üzerinde görüldüğü gibi sınır, enerjiyi en geç başlangıç yörüngesinin üzerinde ya da bu yörünge üzerinde tutar.

\begin{figure}[t]
\centering
\figchain{\linewidth}{\analysisdir/figures/fig2_concepts.pdf}{\analysisdir/figures_conceptual/fig_concepts.pdf}{figures/fig2_concepts.pdf}
\caption{Güvenlik katmanının kavramları. (a) Bir EV'nin kalış süresi boyunca enerjisi; en geç başlangıç yörüngesi $\Bh_i$ ve adım başına $\Delta t\,u_i$, $\Delta t\,l_i$ ve $\Delta t\,\pmn_i$ seçenekleriyle birlikte. (b) İki EV için $-\Crev\le p_1+p_2\le\Ctil$ dilimiyle kesilen $[l,u]$ kutusuna izdüşüm.}
\label{fig:concepts}
\end{figure}

$\Bh_i(k)$, $k=0,\dots,\tau_i(t_i^{\sm{arr}})$ tablosu varışta $\varphi_i$ üzerinde ikiye bölme (bisection) yöntemiyle bir kez hesaplanır; böylece sınırların maliyeti EV ve adım başına $O(1)$ olur.
Saklanan her tablo
\begin{equation}\label{eq:certified}
\Bh(0)\ge\Ereq,\qquad \varphi\big(\Bh(k)\big)\ge\Bh(k-1),\qquad \Bh(k)\le\max\{\Bh(k-1),E^{\min}\},\qquad \Bh(k)\in[E^{\min},E^{\sm{cap}}]
\end{equation}
anlamında sertifikalıdır ve ispatlar yalnızca bu özellikleri kullanır.
Adım cinsinden ölçülen esneklik payı $\chi_i(t)=\big(E_i(t)-\Bh_i(\tau_i(t))\big)/\big(\eta_i^{\sm c}\Pnom_i\Delta t\big)$, ancak ve ancak EV bireysel olarak olurlu olduğunda negatif değildir.

% ---------------------------------------------------------------------
\subsection{İzdüşüm}\label{sec:projection}

$x\in\R^N$, kW cinsinden ham eylem olsun; boş şarj üniteleri için $x_i=l_i=u_i=0$ alınır.
Güvenlik katmanı, sıkılaştırılmış ileri yön limiti olan $\Ctil(t)=C(t)-n(t)\varepsilon_q$ ifadesini kullanır; burada şarj ünitesi başına $\varepsilon_q=0.008$ kW'lık pay, simülatörün eyleyici yuvarlamasını karşılar (Bölüm~\ref{S-app:actuation}).
Olurlu küme $\mathcal K(t)=\{p\in\R^N:\ l\le p\le u,\ -\Crev(t)\le\mathbf 1^\top p\le\Ctil(t)\}$ olup güvenlik katmanı şu Öklid izdüşümüdür:
\begin{equation}\label{eq:proj}
\Pi_t(x)=\arg\min_{p\in\mathcal K(t)}\tfrac12\|p-x\|_2^2 .
\end{equation}

\begin{proposition}[İzdüşüm]\label{prop:projection}
\begin{enumerate}
\item[(i)] $\mathcal K(t)$ kümesi ancak ve ancak $l\le u$, $\sum_il_i\le\Ctil$, $\sum_iu_i\ge-\Crev$ ve $-\Crev\le\Ctil$ ise boş değildir.
\item[(ii)] $\mathcal K(t)\neq\emptyset$ ise \Eq{eq:proj} probleminin tek minimize edicisi $p_i^{*}=\clip(x_i-\lambda^{*},l_i,u_i)$'dir. Kırpılmış nokta toplam limiti zaten sağlıyorsa ortak kaydırma $\lambda^{*}$ sıfırdır; aksi hâlde $\lambda^{*}$, $\lambda\mapsto\sum_i\clip(x_i-\lambda,l_i,u_i)$ fonksiyonunun ihlal edilen limitteki köküdür.
\end{enumerate}
\end{proposition}

İspat Karush--Kuhn--Tucker (KKT) koşullarını doğrulamaktadır; ek materyal ayrıca Jacobi matrisini de vermektedir (Önerme~\ref{S-prop:jacobian}).
\Fig{fig:concepts}(b) üzerinde görüldüğü gibi izdüşüm, tüm serbest koordinatları aynı miktarda kaydırır ve toplam limit sağlanana kadar kırpar.
$2N$ kırılma noktasının $\{x_i-u_i\}\cup\{x_i-l_i\}$ bir kez sıralanması ve bunlar üzerinden bir kez geçilmesi $\lambda^{*}$ değerini verir; dolayısıyla katmanın maliyeti adım başına $O(N\log N)$'dir.
İkiye bölme yöntemine dayalı yedekli bir son denetim, kayan nokta yuvarlamasına karşın toplam limiti korur; Algoritma~\ref{S-alg:safety} tam bir adımı listelemektedir.

Bu cebir yeni değildir: \Eq{eq:proj} sürekli kuadratik sırt çantası problemidir \cite{kiwiel2007breakpoint,wang2022efficient}; sLLF hız kuralı da buna çok yakın bir biçim, yani her EV için tepe hızıyla ölçeklenen ortak bir düzey kullanır \cite{chen2022smoothed}.
Bizim kullanımımız iki bakımdan farklıdır: girdi, boyuttan bağımsız bir politikanın öğrenilmiş şarj ünitesi başına eylemidir; $l_i$ alt sınırı ise şarj gücü azalmasını ve V2G işaret aralığını içeren geriye doğru erişilebilirlikten türetilir.
Ayrıca toplam kısıtı iki yönlüdür ve katman eğitim sırasında da çalışır.

% ---------------------------------------------------------------------
\subsection{Olurluluğun korunması}\label{sec:invariance}

Sezgisel olarak güvenlik katmanı, hiçbir EV'nin gereksinimine tam güçte hâlâ erişilebilen en küçük enerjinin altına düşmesine izin vermez; ortak kaydırma ise istasyonu trafo bandı içinde tutar.
Bu özelliklerden yalnızca ilki sağlanamayabilir; bu durum, trafo her EV'ye alt sınırını aynı anda veremediğinde ortaya çıkar.
$\mathcal K(t)$ boş olmadığında bir adımı olurlu olarak adlandırıyoruz; olursuz bir adımda izdüşümün yerini bir yedek kural alır.
Sonucun koşulları aşağıda toplanmıştır ve bunlardan yalnızca ilk beşi varsayımdır.
Teknik açıdan, sonucun dayandığı tümevarım Chen ve arkadaşlarının esneklik payı argümanını \cite{chen2022smoothed} şarj gücü azalmasına, V2G deşarjına, iki yönlü bir bağlantıya ve keyfî ham eylemlere genişletmektedir.

\begin{center}
\textbf{Teorem~\ref{thm:invariance}'in koşulları}\\[1mm]
\footnotesize
\setlength{\tabcolsep}{3pt}
\begin{tabular}{@{}l>{\raggedright}p{6.4cm}>{\raggedright\arraybackslash}p{6.2cm}@{}}
\toprule
Etiket & Koşul & Bu çalışmadaki durumu \\
\midrule
A1 & $\varphi$ sürekli ve azalmayan & EV2Gym'in şarj gücü azalması için ispatlanmıştır (Lemma~\ref{S-lem:taper}) \\
A2 & ayrılış adımı $d_i$ ve gereksinim $\Ereq_i$ varışta biliniyor & konformal sınır ile gevşetilmiştir (\Sec{sec:conformal}) \\
A3 & tek adım modeli: \Eq{eq:limits} içindeki komutlar için $E_i(t+1)\ge E_i(t)+\Delta t\,g_i(p_i)$ ve $p_i\le\tilde p_i\le p_i+\varepsilon_q$ & EV2Gym'in eyleyici davranışı için ispatlanmıştır (Lemma~\ref{S-lem:a3}); donanımda sağlanamayabilir \\
A4 & olurlu varış: $E_i^{\sm{arr}}\ge\Bh_i(\tau_i(t_i^{\sm{arr}}))$ & talep işleme hattı tarafından uygulanır \\
A6 & şebeke kapasite payı: $C(t)\ge N\varepsilon_q$ ve $\Crev(t)\ge0$ & dışsaldır; yedek kural sayımlarıyla denetlenir \\
\midrule
A5 & olurlu adım: $\sum_il_i\le\Ctil$, $\sum_iu_i\ge-\Crev$, $\Ctil\ge-\Crev$ & kapalı döngüde ortaya çıkan bir olaydır; (a) ve (b) tarafından varsayılmaz; ölçülür \\
\bottomrule
\end{tabular}
\end{center}

A6, politikaya bağlı olmaması için $n(t)$ yerine şarj ünitesi sayısı $N$'yi kullanır ve senaryolarımızda epizot başına $2\times10^{-9}$'un altında bir olasılıkla sağlanamaz (Bölüm~\ref{S-app:a6}).

\begin{theorem}[Olurluluğun korunması]\label{thm:invariance}
Eğitim ya da kullanım sırasında, sonlu ham eylemlerden oluşan keyfî bir dizi için komut edilen güçler olurlu adımlarda $p(t)=\Pi_t(x(t))$, diğer adımlarda ise yedek kuralın çıktısı olsun.
\begin{enumerate}
\item[(a)] A3 ve A6 altında, her adımda gerçekleşen güçler $-\Crev(t)\le\sum_i\tilde p_i(t)\le C(t)$ trafo limitini, komut edilen güçler ise \Eq{eq:limits} ile verilen fiziksel limitleri sağlar.
\item[(b)] A1--A4 altında, kalış süresi boyunca $p_i(t)\ge l_i(t)$ olan her EV $i$, $E_i(d_i+1)\ge\Ereq_i$ koşulunu sağlar; özellikle kalış süresi boyunca olursuz adımla karşılaşmayan her EV bu koşulu sağlar.
\item[(c)] Ek olarak her adım olurlu ise her EV ayrılış gereksinimini karşılar.
\end{enumerate}
\end{theorem}

\begin{proof}[İspat taslağı]
Olurlu adımlarda Önerme~\ref{prop:projection}(i) $p(t)\in\mathcal K(t)$ sonucunu verir ve A3, sıkılaştırılmış bandı gerçekleşen güçlere aktarır.
A6 altında yedek kural yalnızca aşağıda açıklanan gerekli güç fazlası durumuyla karşılaşabilir ve bu durumda çıktısı, komut edilen toplamı $\Ctil(t)$ değerinde tutar.
(b) için tek adımlık bir lemma (Lemma~\ref{S-lem:onestep}), $E_i(t)\ge\Bh_i(\tau_i(t))$ değişmezini A4'ün sağladığı varıştan itibaren $p_i\ge l_i$ olan her adım boyunca taşır; son adımda bu değişmez $E_i(d_i+1)\ge\Bh_i(0)\ge\Ereq_i$ sonucunu verir.
(c) şıkkı (b)'den çıkar; tam ispat Bölüm~\ref{S-app:theorem} içinde verilmektedir.
\end{proof}

Trafoya ilişkin kısım ne politikaya ne de A1 ya da A5'e bağlıdır; ayrılışa ilişkin kısım ise araç başına bir ifadedir.
A5, ortak bir plan mevcut olduğunda bile sağlanamayabilir.
10 kW'lık bir limit altında her biri 1 saatlik iki adım içinde 10 kWh'e ihtiyaç duyan iki EV, kapasiteyi eşit biçimde paylaşabilir.
Ancak birinci adımda boşta bekleyen bir politika, ikinci adımda $\sum_il_i=20$ kW bırakır.
Bu nedenle EV'lerin olursuz adımlarla ne sıklıkta karşılaştığı ölçülmektedir.

% ---------------------------------------------------------------------
\subsection{Yedek kural}\label{sec:fallback}

Olursuz bir adımda yedek kural, izdüşümün yerine en az esneklik payı önce ilkesine dayanan deterministik bir paylaştırma koyar.
İlgili durumda gerekli güç $\sum_il_i$, $\Ctil$ değerini aşar.
Bu durumda her EV fiziksel alt limiti $\pmn_i$'den başlar ve $\Ctil$ değerine kadar kalan bütçe, esneklik payı artan sırada EV'lere, her biri en fazla $l_i$'ye kadar olmak üzere dağıtılır.
Bu çıktı trafo limitini korur ve hâlâ $l_i$ alan her olurlu EV olurlu kalır (Lemma~\ref{S-lem:fallback}).
Üç ek durum çelişkili limitleri, aşırı PV ihracatını ve negatif bir bütçeyi kapsar; A6 bu üçünü de dışlar (Bölüm~\ref{S-app:fallback}).
$l_i$'den daha az alan EV'ler için bir epizottaki toplam ayrılış eksikliği, bu adımlarda biriken açık ile önsel olarak sınırlanır (Lemma~\ref{S-lem:deficit}).
Talep işleme hattındaki kabul kontrolü olursuz adımları daha seyrek hâle getirir; ancak tahminlere dayanır ve Teorem~\ref{thm:invariance} kapsamına girmez.

% ---------------------------------------------------------------------
\subsection{Ayrılış adımı için konformal alt sınır}\label{sec:conformal}

Ayrılış adımı açısından A2'yi gevşetmek için tek yönlü, kantil regresyonu biçiminde \cite{romano2019conformalized} bölünmüş konformal tahmin (CP) \cite{papadopoulos2002inductive,vovk2022algorithmic,lei2018distributionfree,angelopoulos2023conformal} kullanılmaktadır.
Bir gradyan artırma (gradient boosting) regresörü $\hat q$, $d_i$'nin alt kantilini varış adımından, haftanın gününden, varış enerjisinden, talep edilen enerjiden ve EV modelinden kestirir.
Kalibrasyon, bağımsız tohumlarla (seed) üretilmiş $n_{\sm{cal}}$ epizodun her birinden düzgün dağılımla bir oturum seçer ve $e_j=\hat q(\mathbf w_j)-d_j$ skorlarını kullanır.
$\hat Q_{1-\alpha}$, $\lceil(n_{\sm{cal}}+1)(1-\alpha)\rceil$'inci en küçük skor olmak üzere $\hat d_i^{\sm{lo}}=\max\{t_i^{\sm{arr}},\lceil\hat q(\mathbf w_i)-\hat Q_{1-\alpha}\rceil\}$ sınırı; sınırlarda, talep işleme hattında ve özniteliklerde $d_i$'nin yerini alır ve $\hat d_i^{\sm{lo}}$ adımından sonra katman tek bir kalan adım kullanır.

\begin{proposition}[Konformal ayrılış garantisi]\label{prop:conformal}
Kalibrasyon oturumları ile test oturumu değiştirilebilir (exchangeable) olsun.
Bu durumda $\mathbb P(d_i\ge\hat d_i^{\sm{lo}})\ge1-\alpha$ olur ve $d_i$ yerine $\hat d_i^{\sm{lo}}$ alınarak A1, A3 ve A4 altında $\mathbb P\big(E_i(d_i+1)\ge\Ereq_i\big)\ge1-\alpha-\mathbb P\big(\exists\,t\le d_i:\ p_i(t)<l_i(t)\big)$ sağlanır.
\end{proposition}

Burada $\Ereq_i$, $\hat d_i^{\sm{lo}}$ ile yeniden sertifikalanmış gereksinimdir ve olasılık, kalibrasyon verisi ile bir test epizodundan düzgün dağılımla seçilen bir EV üzerinden alınır; dolayısıyla garanti marjinaldir.
Kalibrasyon her test yapılandırmasından veri gerektirdiğinden konformal varyant sıfır-atış değildir.

% ---------------------------------------------------------------------
\subsection{Eğitim algoritması}\label{sec:training}

Politika, güvenlik katmanı, yedek kural ve eyleyici davranışı ortamın parçası olarak ele alınarak yakınsal politika optimizasyonu (PPO) \cite{schulman2017proximal} ve genelleştirilmiş avantaj kestirimi \cite{schulman2016highdimensional} ile eğitilir.
Bu nedenle PPO, ham Gauss örneğinin olabilirlik oranını kullanır.
Eski parametrelerdeki PPO vekil amaç gradyanına eşit olan temel (vanilla) politika gradyanı kestiricisi, bu durumda izdüşümü alınmış eylemlerin getirisi için yansızdır.
Stokastik politikalar için ortamın içindeki ve politikanın içindeki izdüşüm bu anlamda denktir \cite{markgraf2025safe}; bu tercih aktif bir izdüşümün sıfır gradyanlarından da kaçınır \cite{kasaura2023benchmarking}.
$u_i-l_i\le10^{-9}$ kW olan belirteçler hiçbir komut gücünü etkileyemeyeceğinden log-olasılıktan ve entropiden çıkarılır.
Yedek kural adımları aktör kaybında maskelenir, ancak kritik için tutulur.
Birçok ham eylem aynı izdüşüm noktasına eşlendiğinden olurlu adımlarda $\beta=0.01$ ağırlıklı bir düzenlileştirici, politika ortalamasını $\mathcal K(t)$'ye yakın tutar:
\begin{equation}\label{eq:reg}
\mathcal L_{\sm{reg}}(\theta)=\beta\,\mathbb E_{s}\Big[\frac1{n(t)}\sum_{i\in\mathcal N(t)}\Big(\frac{\bar x_i-\Pi_t(\bar x)_i}{P^{\sm{ref}}}\Big)^{2}\Big],\qquad \bar x=P^{\sm{ref}}\mu_\theta(s).
\end{equation}
Moreau özdeşliğine göre $\tfrac12\mathrm{dist}^2(x,\mathcal K)$'nin gradyanı $x-\Pi_t(x)$'tir; bu nedenle $\Pi_t(\bar x)$ gradyan durdurma (stop-gradient) altında hesaplanır \cite{chen2021enforcing,markgraf2025safe}.
Her veri toplama turu (rollout) $N=10$ olan sekiz ve $N=20$ olan sekiz ortam çalıştırır; tam eğitim döngüsü Algoritma~\ref{S-alg:ppo} içinde verilmektedir.
